\pdfoutput=1
\documentclass[11pt]{amsart}
\ifdefined\pdftrailerid\pdftrailerid{}\fi
\usepackage[margin=1.2in]{geometry}
\usepackage{amsmath,amssymb,amsthm,mathtools}
\usepackage[colorlinks=true,linkcolor=blue,citecolor=blue,urlcolor=blue]{hyperref}
\hypersetup{pdftitle={The quadratic and cubic coefficients of circle exchange for balanced local curves},pdfauthor={Bernd Johannes Wuebben}}
\setlength{\emergencystretch}{2em}
\newtheorem{theorem}{Theorem}[section]
\newtheorem{proposition}[theorem]{Proposition}
\newtheorem{lemma}[theorem]{Lemma}
\newtheorem{corollary}[theorem]{Corollary}
\theoremstyle{definition}
\newtheorem{definition}[theorem]{Definition}
\newtheorem{example}[theorem]{Example}
\theoremstyle{remark}\newtheorem{remark}[theorem]{Remark}
\DeclareMathOperator{\Pic}{Pic}
\DeclareMathOperator{\Sym}{Sym}
\DeclareMathOperator{\CT}{CT}
\DeclareMathOperator{\Res}{Res}
\newcommand{\OO}{\mathcal O}
\newcommand{\ZZ}{\mathbb Z}
\newcommand{\QQ}{\mathbb Q}
\newcommand{\CC}{\mathbb C}
\title[Quadratic and cubic circle exchange]{The quadratic and cubic coefficients of circle exchange for balanced local curves}
\author{Bernd Johannes Wuebben}
\subjclass[2020]{14N35 (primary), 14C35, 05E05 (secondary)}
\keywords{Stable pairs, local curves, equivariant $K$-theory, Macdonald operators, Picard varieties}
\date{November 27, 2026}
\begin{document}
\begin{abstract}
Exchanging two circle actions on a Calabi--Yau fivefold gives stable-pair indices on a balanced local curve and on the product of its base curve with an affine plane. After a common determinant normalization, we prove that the two modified indices agree through cubic order in the exchanged counting variable, in integral $K$-theory over the Picard variety parametrizing one normal summand. The comparison holds for every smooth projective complex curve of positive genus and every positive curve degree, at every Euler characteristic and arbitrary volume character. On the formal disk we identify the quadratic and cubic coefficients separately for each Young diagram, with the cotangent character also free. One presentation is a Macdonald expectation and the other is a finite sum over reverse plane partitions. At quadratic order the residues from interchanging currents give the two orientations of a domino. At cubic order an exact polynomial interpolation proves that two adjacent diagonal parameters can be interchanged in the sum with currents grouped by sign. Restoring the original order produces an additional pole whose contributions from the ordinary currents and the second Adams operation cancel identically.
\end{abstract}
\maketitle
\section{Introduction}\label{sec:intro}

A choice of circle action on a Calabi--Yau fivefold can determine both a fixed threefold and the variable that counts the Euler characteristic of its stable pairs. Changing that circle can therefore exchange a counting variable with a character of the normal torus. Nekrasov and Okounkov formulated this principle in their proposed correspondence between membranes on fivefolds and sheaves on threefolds \cite[\S\S 1--2]{NO}. For a balanced local curve the two fixed threefolds have different deformation theories. One is a local Calabi--Yau threefold varying in a Picard family. The other is the product of the curve with an affine plane. We compare their family indices through degree three in the exchanged counting variable, with the volume character free.

\subsection{The two threefolds and their gradings}
Let $C$ be a smooth projective complex curve of genus $g\geq 1$, let $b=g-1$, and put $B=\Pic^b(C)$. Choose a point of $C$ and let $L$ denote the Poincar\'e line bundle on $C\times B$ trivialized over that point. We write $K$ for the pullback of $K_C$. Both summands of $L\oplus K L^{-1}$ have degree $b$ on $C$, which is the balance condition. On the total space of
\[
         tL\;\oplus\;(\kappa/t)K L^{-1}
             \;\oplus\;q\OO\;\oplus\;(\kappa^{-1}/q)\OO
\]
the products of the two pairs of fiber characters are inverse to one another. The subgroups act on the four summands by $(1,1,z,z^{-1})$ and $(z,z^{-1},1,1)$, respectively. The first fixes the total space of the first pair of lines. The second fixes $C\times\mathbb A^2$ over $B$. In both cases the external directions enter the determinant of the stable-pair index. On the product threefold this is how the index retains the dependence on $L$.

A stable pair consists of a pure one-dimensional sheaf $\mathcal F$ of proper support and a section with zero-dimensional cokernel. Its curve degree $m$ is the generic length over $C$. Denote the two oriented modified stable-pair indices pushed forward to $B$ by $\mathcal I_m^{\rm orig}$ and $\mathcal I_m^{\rm rot}$. We use the ordered external lines above and the determinant square root of \cite[Lemma~2.1]{W1} on the cover $\sigma^2=\kappa$. The first index counts $\mathcal F$ with $(-q)^{\chi(\mathcal F)}$. The second counts with $(-t)^k$, where $k=\chi(\mathcal F)+mb$. On a torus-fixed pair on the product threefold, $k$ is the sum of the degrees of the effective divisors that modify its monomial summands. Two units can belong to the same summand, or to different summands with incident support. The quadratic and cubic coefficients are the coefficients of $t^2$ and $t^3$.

\begin{theorem}\label{thm:familyintro}
Let $m\ge1$ and $D_m=bm(m+1)/2$. In the common expansion at $q=t=0$,
\[
 \sigma^{mb}\mathcal I_m^{\rm orig}\equiv
 \sigma^{mb}\mathcal I_m^{\rm rot}\pmod {t^4}
 \qquad\text{in }q^{-D_m}K^0(B)[\kappa^{\pm1}][[q,t]].
\]
The congruence means equality of the coefficients of $t^k$ for $0\leq k\leq 3$ at each Euler characteristic.
\end{theorem}

The coefficient ring is the algebraic $K$-group of the entire Picard variety. Passing to the rank on a fiber can lose the class under comparison. For example, in genus two and curve degree two the coefficient of $(1-\kappa)q^{-3}t^2$ of the unnormalized index, expanded about $\sigma=1$, is $-6\psi^2 R\pi_*L$, where $\pi:C\times B\to B$ and $\psi^2$ is the second Adams operation \cite[Theorem~1.1]{W1}. This class has rank zero and first Chern character $12\theta$ in rational cohomology, for $\theta$ the principal polarization. Theorem~\ref{thm:familyintro} retains the full volume character at this grading, at all higher Euler powers of the same curve degree, and in every positive curve degree.

\subsection{A local identity for each partition}
The two full-torus localization formulas are indexed by partitions $\lambda$ of $m$. Write $h_u$ and $\ell_u$ for the hook and leg lengths at a box $u$, and put $H_\lambda=\sum_{u\in\lambda}h_u$. The lowest Euler power for this partition is $-bH_\lambda$. On a formal disk the excess above that minimum has counting variable $x$ and the number of exchanged marks has variable $y$. Write $p$ for the cotangent character of the disk, $c=\kappa^{-1}$ and $s=p/c$. Let $A_{\lambda;c,s}(x,y)$ be the stable-pair vertex of equation~\eqref{eq:A}, including its hook product on each coordinate axis. Exchanging the external circles sends $(x,y,c,s)$ to $(y,x,s,c)$.

\begin{theorem}\label{thm:diskintro}
For every Young diagram $\lambda$ and every $n,k\ge0$ with $\min(n,k)\le3$,
\[
 [x^ny^k]A_{\lambda;c,s}(x,y)
       =[x^ny^k]A_{\lambda;s,c}(y,x)
       \qquad\text{in }\QQ(p,c).
\]
For fixed $k\le3$ the two coefficients of $y^k$, divided by the common coefficient of $y^0$, agree as rational functions of $p,c,x$.
\end{theorem}

Equality along the axes follows from the hook product of Dinkins and Smirnov \cite{DS}. The first exchanged coefficient was computed in \cite[Theorem~1.2 and \S6]{W1}. The same paper proved agreement to first order about volume one in both gradings. At volume one the earlier comparison holds on the common space of cycles \cite[Theorem~1.1]{W0}. Theorem~\ref{thm:familyintro} concerns the pushforwards to $B$. Its diagramwise form compares the localized contributions at $q^{-bH_\lambda+n}t^k$ when $\min(n,k)\le3$. Since the minima depend on $\lambda$, the integral degree sum is taken at a common absolute Euler power. The local identity retains the independent character $p$ throughout the disk calculation.

\subsection{How the terms match}
In the rotated vertex, extracting $y^k$ leaves a finite sum over weakly increasing arrays of nonnegative heights on $\lambda$ with total height $k$. For $k=2$ there are three possibilities. A corner can be raised twice, two corners can each be raised once, or the two raised boxes can form a removable domino. In the original vertex $x$ records the total height, which is unbounded even at $y^2$. Extracting $y^2$ there inserts a symmetric function of the diagonal heights into the Macdonald measure of the partition. Thus the two sides of Theorem~\ref{thm:diskintro} have summation indices of different sizes.

The Macdonald operators write the insertion as a sum of products acting at successive diagonals of $\lambda$. Each diagonal lies between two steps on its southeast boundary. A positive current can be evaluated on the downward steps to its right and a negative current on the horizontal steps to its left. To do this simultaneously for several insertions one must change their order. In the quadratic term the change has two residues. After summing the locations of the insertions, their contributions are exactly the two orientations of a domino. The terms already in the desired order give two separate corners and a corner of height two. The second Adams operation supplies the coefficient for selecting a step twice.

To sum the possible insertion positions at cubic order, we first group the currents by sign. The resulting sum is invariant under permuting the parameters attached to two adjacent diagonals. Moving the parameter of a selected downward step to the first position forces a corresponding horizontal step at the same corner. Repeating this argument identifies all the terms in the grouped sum. We prove the permutation identity for abstract boundary weights: after its denominators are cleared, the degrees in the two Macdonald parameters are at most 17 and 29. Exact interpolation on an $18\times 30$ grid, together with bounds in the other variables, proves the identity for independent characters. The calculation involves at most three labels of each sign, however many boxes the original diagram has.

Changing the currents back to their original order requires a separate contour calculation. With three marks, a pair of currents can interact with the field remaining after the first residue. In addition to a shift by $p$, the remaining integral has a pole at a boundary weight divided by $c$. Its contribution appears both in the product of three ordinary insertions and in the term with a second Adams operation. Each contribution can be nonzero even on a diagram consisting of a single row of two boxes. The two cancel after summing the positions on the boundary. The surviving terms include three successive boxes, a domino and another mark, and stacks at removable corners. The proof requires both the cancellation and the enumeration of the possible contour residues.

Section~\ref{sec:disk} sets the conventions for the two disk calculations. Sections~\ref{sec:quadratic} and~\ref{sec:cubic} compare their coefficients of orders two and three. Section~\ref{sec:family} descends the identity across divisors of coinciding marks to the Picard variety. At fourth order two second Adams insertions can cross. The comparison of their residues for arbitrary diagrams lies beyond the results proved here.

\section{The disk vertex and the diagonal operators}\label{sec:disk}

A box of $\lambda$ has coordinates $(i,j)$, starting with $(0,0)$ at the northwest corner, with $i$ increasing to the right and $j$ downwards. The content is $i-j$. A reverse plane partition $d$ is a nonnegative height $d_{ij}$ at each box, weakly increasing to the right and down. Its size is $|d|=\sum d_{ij}$. For a character $U$, the dual $U^*$ inverts each weight. On a finite virtual sum $U=\sum m_w w$ put
\[
 R_a(U)=\prod_w\left(\frac{1-aw}{1-w}\right)^{m_w},
 \qquad F_a(z)=\sum_{j\ge0}\frac{(a;p)_j}{(p;p)_j}z^j,
 \qquad (a;p)_j=\prod_{v=0}^{j-1}(1-ap^v).
\]
We use the same notation $R_a(z)$ for a single weight.

\subsection{The two expansions}
Set $E_d=\sum_{(i,j)\in\lambda} y^{j-i}s^{-j}p^{-d_{ij}}$ and
\[
       H_y(E)=\frac{E-(1-y)EE^*}{1-p}.
\]
The change $H_y(E_d)-H_y(E_0)$ has finite Laurent support. The common convention for the disk vertex is
\begin{equation}\label{eq:A}
 A_{\lambda;c,s}(x,y)=
 \prod_{u\in\lambda}F_c(y^{h_u}s^{-\ell_u})
 \sum_d x^{|d|}c^{-|d|}
 R_c\big(H_y(E_d)-H_y(E_0)\big).
\end{equation}
At fixed height, collect equal Laurent monomials before evaluating $R_c$. Reflect negative powers of $y$ before Taylor expansion at the origin. A weight-$w$ summand of a monomial pair on the disk is a copy of the divisor line with fiber $wp^{-d}$ at the origin. The two coordinate multiplication arrows give the height inequalities, including both arrows into a box with two predecessors. The formula uses the fixed virtual structure sheaf and the relative determinant at these divisors \cite[Theorem~6.1]{Mon}, \cite[\S5]{W1}. The factor $c^{-|d|}$ includes the ordered external determinant.

Write $Z=A_{\lambda;c,s}(x,0)=A_{\lambda;s,c}(0,x)$ for the axis series of \eqref{eq:family-axis} and set
\[
 T_k=Z^{-1}[y^k]A_{\lambda;c,s}(x,y),\qquad
 U_k=Z^{-1}[y^k]A_{\lambda;s,c}(y,x).
\]
The sums for $T_k$ extend over all original heights. On the other side put $v=c/x$, $w_{ij}=x^{-i}v^{-j}$, and $E=\sum_{\lambda}w_{ij}$. For the exchanged heights $e$, define
\[
 Q_e=\sum_{(i,j)\in\lambda}w_{ij}
            \sum_{0\le a<e_{ij}}p^{-a},\quad
 B_\lambda=\sum_{\rm removable}w_{ij}-c\sum_{\rm addable}w_{ij},
 \quad D=(p^{-1}-1)(1-x).
\]
A removable box is one whose deletion leaves a partition, and an addable box is one whose addition does so. Then
\begin{equation}\label{eq:W}
 U_k=\sum_{|e|=k}W(e),\qquad
 W(e)=R_s\big(B_\lambda Q_e^*+D Q_eQ_e^*\big).
\end{equation}
To obtain this form, substitute $E+(p^{-1}-1)Q_e$ into $H_x$ and subtract $H_x(E)$. The difference is
\[
 p^{-1}Q_e+(1-x)(EQ_e^*-p^{-1}Q_eE^*)+D Q_eQ_e^* .
\]
Reflect the first and third linear terms using $R_s(z)R_s(1/(sz))=s$, and use $(1-x)(1-v)E-c=B_\lambda$, obtained by summing the rows of the diagram. The constant from reflection cancels $s^{-|e|}$. At a monomial of weight $1$ or $1/s$ the reflected product means its value in the original finite polarization in \eqref{eq:A}, with determinants and equal factors combined before taking the value. This convention also applies to the products of adjacent box weights below.

\subsection{The boundary measure}
Put $J_i=\max(0,-i)$, let $v_i$ be the number of boxes of content $i$, and set $L_i=J_i+v_i$. The vertices $(i+L_i,L_i)$ form the southeast boundary from $(0,\operatorname{height}\lambda)$ to $(\operatorname{width}\lambda,0)$. Write $Z_i=x^ic^{L_i}$. Along a horizontal step (H) the ratio $Z_{i+1}/Z_i$ equals $x$, and along a downward step (D) it equals $x/c$. A slice at $i$ is the content diagonal immediately before step $i$. Its alphabet of weights for the height $d$ is
\[
 X_i=\sum_{(i+j,j)\in\lambda}s^j p^{d_{i+j,j}},
 \qquad V=\sum_i y^i X_i .
\]
Write $\langle\,\cdot\,\rangle$ for the expectation in the height measure of \eqref{eq:A} at $y=0$, divided by $Z$. The termwise interlacing expansion identifies this measure with the Macdonald process of parameters $(p,s)$. For a horizontal step use the specialization $1/Z_\beta$ and for a downward step $Z_\alpha$. A horizontal step followed by a downward step contributes $F_s(Z_\alpha/Z_\beta)$ to the vacuum matrix element. Multiplication over such pairs gives $Z$ by the hook-product formula \cite[\S3, Propositions~6--7]{DS}, \cite[Lemma~5.1]{W1}.

Put $C_h=[y^h]\{sV+V^*-(1-y)(1-s/y)VV^*\}$ for $h\ge1$ and $l_j=(1-c^j)/(1-p^j)$, with $l=l_1$. Extracting the positive axial weights in \eqref{eq:A} gives
\begin{equation}\label{eq:desc}
 \sum_{k\ge0}y^k T_k
 =\left\langle\exp\left(\sum_{h,j\ge1}
             \frac{l_j}{j}\,y^{hj}\psi^j C_h\right)\right\rangle .
\end{equation}
Here $\psi^j$ raises every character in $C_h$ to the $j$th power before taking expectations. Indeed the positive part of the finite-height ratio is the symmetric algebra on $(1-c)/(1-p)$ times the change in the positive weights of the displayed character from the empty heights. At empty heights those weights are $\sum_u y^{h_u}s^{-\ell_u}$, supplied by the prefactor of \eqref{eq:A}. The remaining degree-zero factors give the interlacing measure \cite[\S5.2--5.3]{W1}. The exponential in \eqref{eq:desc} takes place over the disk character field.

\subsection{Ordering the insertions}
Write
\[
 P_i=(s-1)X_i+s^{J_i},\quad N_i=(s^{-1}-1)X_i^*+s^{-J_i},
 \quad \gamma=-\frac{s}{(s-1)^2}.
\]
These are the eigenvalues of the scaled first Macdonald operators of the slice and of its dual. The differences give
\begin{equation}\label{eq:Ch}
 C_h=\gamma\sum_i (N_{i+1-h}-sN_{i-h})(P_i-P_{i+1}).
\end{equation}
For $E_j'=\sum_iX_{i+j}X_i^*$ the quadratic part is
$E_{h-1}'-(1+s)E_h'+sE_{h+1}'$. The mixed terms with $s^{\pm J_i}$ telescope to $sX_h+X_{-h}^*$, and the two constant differences have disjoint supports. This proves \eqref{eq:Ch} even at the first and last occupied diagonals. We call the summation index $i$ in $C_1$ its center. A cut is the position immediately before an actual boundary step. A term centered at $i$ chooses its positive field at cut $i$ or $i+1$ and its negative field at cut $i$ or $i-1$. Thus several labeled insertions can share a cut even when their centers differ.

For use in products, insert the currents in increasing slice order. At a common slice the zero modes are diagonal in the same Macdonald basis, so we may put negatives before positives there. With $b_n$ the power-sum coordinate, the field with positive zero mode has creation and annihilation factors
\[
 \exp\!\sum_{n>0}(1-s^{-n})b_n z^n/n,\qquad
 \exp\!\sum_{n>0}(p^n-1)z^{-n}\partial_{b_n}.
\]
For the negative field invert $p$ and $s$. The two fields are scaled by $s^{L_i}$ and $s^{-L_i}$, respectively, to give the eigenvalues above \cite[Propositions~2.6--2.7]{Ko}. Moving an annihilation factor past a creation factor gives the kernel
\begin{equation}\label{eq:K}
 K_{a,b}(h)=\frac{(1-h)(1-ah/b)}{(1-ah)(1-h/b)},\qquad
 h=\frac{\text{later spectral variable}}{\text{earlier spectral variable}} .
\end{equation}
Here $a=p$ or $p^{-1}$ as the earlier field is positive or negative, and $b=s$ or $s^{-1}$ as the later field is positive or negative. The expansion is in nonnegative powers of $h$. Commutation with the boundary steps leaves the following one-variable factors:
\begin{equation}\label{eq:onebody}
\begin{aligned}
 P_i(z)&=\prod_{\beta<i,\mathrm H}\frac{1-z/(s Z_\beta)}{1-z/Z_\beta}
          \prod_{\alpha\ge i,\mathrm D}\frac{1-z/Z_\alpha}{1-z/(sZ_\alpha)},\\
 N_i(w)&=s^i\prod_{\beta<i,\mathrm H}\frac{1-Z_\beta/(sw)}{1-Z_\beta/w}
          \prod_{\alpha\ge i,\mathrm D}\frac{1-Z_\alpha/(cw)}{1-Z_\alpha/(pw)} .
\end{aligned}
\end{equation}
Their boundary contributions on shrinking $z$ and expanding $w$ are
\[
 A_i^\alpha=\Res_{z=sZ_\alpha}P_i(z)\,dz/z,\qquad
 B_i^\beta=-\Res_{w=Z_\beta}N_i(w)\,dw/w .
\]
An unavailable step has coefficient zero. The constants at the two ends, $1$ and $s^i$, cancel in the forward and backward differences in \eqref{eq:Ch}. One can first take the downward roots small, the horizontal roots large, and nested spectral radii between them. Calculations made in separate compatible annuli are continued as rational identities before specializing $Z_i=x^ic^{L_i}$.

Higher Adams insertions act by the Newton polynomials in the elementary Macdonald operators. For clarity, in an alphabet $a_1,\dots,a_m$ their order-$h$ elementary operator with parameters $(u,v')$ is
\[
 (v')^{\binom h2}
 \sum_{|I|=h}\prod_{i\in I,j\notin I}\frac{v'a_i-a_j}{a_i-a_j}\,T_{u,I},
\]
where $T_{u,I}$ multiplies the selected variables by $u$. Its eigenvalue on a partition $\mu$ is $e_h((u^{\mu_j}(v')^{m-j})_{1\le j\le m})$ \cite[\S2.2]{BC}. The preceding horizontal alphabet may be used from the left and the following downward alphabet from the right. Inverting both parameters preserves the monic Macdonald basis by homogeneity of its scalar product. Passing to normalized moments conjugates a shift by the corresponding vacuum factors $F_s$; their shifted ratios are finite products. In particular an earlier shift acts on the coefficients of every later operator. Equivalently, the ordered constant term contains one copy of \eqref{eq:K} for each ordered pair of fields. Sections~\ref{sec:quadratic} and~\ref{sec:cubic} use these two descriptions to compare the complete expectations with \eqref{eq:W}.

\section{The quadratic disk coefficient}\label{sec:quadratic}

\begin{theorem}\label{thm:quadraticdisk}
For every Young diagram the two disk vertices, divided by their common
value at $y=0$, have the same coefficient of $y^2$ in $\QQ(p,c,x)$.
\end{theorem}

A reverse plane partition of size two has two distinct maximal boxes of
height one, a maximal box of height two, or a removable domino of height one.
We recover the first two terms by collecting the currents of each sign.
Interchanging two currents of opposite signs gives the domino term.

\subsection{Interchanging two currents}
Write $R=R_s$, and use the residues of \eqref{eq:onebody}. Thus $A_i^\alpha$ is the signed inner residue of $P_i(z)$ at
$sZ_\alpha$ and $B_i^\beta$ is the signed outer residue of $N_i(w)$ at
$Z_\beta$. Thus $\alpha$ is a downward step at or after $i$ and $\beta$
a horizontal step before $i$. A residue outside its indicated range is
zero. Use the differences $\Delta A_i=A_i-A_{i+1}$,
$\nabla B_i=B_i-sB_{i-1}$. Set
\[
\begin{aligned}
 r_z(X)&=\frac{X-z}{X-sz},& S(u)&=\frac{R(u)}{R(cu)},\\
 M(u,n)&=\frac{(u;p)_n}{(su;p)_n},&
 Q_{m,n}(u)&=\frac{M(u,m+n)}{M(u,m)M(u,n)} .
\end{aligned}
\]
Here $(u;p)_{-n}=1/(p^{-n}u;p)_n$ for $n>0$. We use the transport laws
\begin{equation}\label{eq:qtrans}
 A_{i+1}^\alpha=r_{Z_\alpha}(Z_i)A_i^\alpha,\qquad
 B_{i-1}^\beta=r_{xZ_\beta}(Z_i)B_i^\beta .
\end{equation}
For the first, $P_{i+1}(z)/P_i(z)=(1-z/(sZ_i))/(1-z/Z_i)$ at either
type of step. For the second, $N_i(w)/N_{i-1}(w)$ equals
$R(w/Z_{i-1})$ at a preceding horizontal step and $R(cw/Z_{i-1})$
at a preceding downward step. The increments of $Z_i$ are $x$ and $x/c$,
respectively. Both relations hold at the physical endpoints by clearing
denominators.

Let $G_{11}$ be the following evaluation of $C_1^2$. Expand the slice
differences of each factor. Evaluate the two negative fields on their
horizontal prefixes in increasing slice order, and the two positive fields
on their downward suffixes in decreasing slice order. On ties use center
order and then label order in the chronological word, reversing that order
for the suffix action. The finite Macdonald difference operators retain all
prior shifts of their own alphabet. For each branch with horizontal
multiplicities $m_\beta$ and downward multiplicities $n_\alpha$, insert
\[
                 \prod_{\beta,\alpha}
                 Q_{-m_\beta,n_\alpha}(Z_\alpha/Z_\beta).
\]
This prescription includes all pairs of selected roots. Constants in a
negative or positive field cancel in its own slice difference: their
values are $s^i$ on the prefix and $1$ on the suffix, and the constant
operator shifts no variable. Thus each nonzero branch selects two roots
of each sign.

For $k\le j$, let $E_{NP}$ evaluate the pair $N_j,P_k$ by the indicated
prefix and suffix with its $Q$ factor, including constant choices.
Let $E_{PN}$ instead apply $P_k$ on the prefix and $N_j$ on the suffix,
with their positive and negative shifts. The latter evaluates the
chronological pair. Put $\omega=(1-s)(1-p^{-1})/(1-s/p)$.
\begin{lemma}\label{lem:quadraticinterchange}
For independent boundary roots,
\begin{equation}\label{eq:qmiddle}
 E_{NP}-E_{PN}=
 \omega(1-s)\sum_{k\le h<j}
 \prod_{\substack{k\le r<j\\r\ne h}}R(Z_h/Z_r)\,(\Phi_{h,j}-1),
\end{equation}
where at a horizontal step $\Phi_{h,j}=\prod_{d\ge j,\,D}S(Z_h/Z_d)$,
and at a downward step $\Phi_{h,j}=\prod_{b<j,\,H}S(Z_h/(cZ_b))$.
Here and below $H,D$ restrict a step index to the indicated type.
\end{lemma}
\begin{proof}
Put $f(u)=u/(1-u)$. Direct division gives
$Q_{-1,1}(u)-1=\omega(f(u)-f(su/p))$.
Place all downward roots sufficiently close to zero and all horizontal
roots sufficiently far from zero for one annulus to separate their poles,
also after multiplication by $s$ and $p$. Expand
\[
       \omega\CT_w\big[(P_k(sw)-P_k(pw))N_j(w)\big].
\]
The constant modes cancel. Pairing negative powers of $P_k$ from an inner
downward pole with positive powers of $N_j$ from an outer horizontal
pole gives the product of their signed residues times $Q_{-1,1}-1$.
In the opposite pairing the inner pole of $N_j$ is at $Z_d/p$ and the
outer pole of $P_k$ is at $Z_b$. The geometric sums give
$\omega(f(sZ_d/(pZ_b))-f(Z_d/Z_b))$, with the opposite sign.
The products of the separate means in the two evaluations agree. Thus
the displayed constant term is their difference. Rational continuation
allows arbitrary independent roots, even when a horizontal step in
$[k,j)$ follows a downward step in that interval.

At the cut $j$ the one-body products simplify to
\[
 P_j(sw)N_j(w)=\prod_{d\ge j,\,D}S(w/Z_d),\qquad
 P_j(pw)N_j(w)=\prod_{b<j,\,H}S(w/Z_b).
\]
Changing $j$ to $k$ in $P_j(z)$ multiplies it by
$\prod_{k\le r<j}(1-z/Z_r)/(1-z/(sZ_r))$.
Enlarge the first contour. The value at infinity is $s^{j-k}$, and
the outer residues occur at $Z_h$ for horizontal $h$ in the interval,
with coefficient $(1-s)\prod_{r\ne h}R(Z_h/Z_r)\Phi_{h,j}$.
Shrink the second contour. The value at zero is one and the inner
residues at $Z_h/c$ for downward $h$ have coefficient
$-(1-s)\prod_{r\ne h}R(Z_h/Z_r)\Phi_{h,j}$.
Finally $\sum_h\prod_{r\ne h}R(Z_h/Z_r)=(1-s^{j-k})/(1-s)$,
by partial fractions of $\prod_{k\le r<j}R(w/Z_r)$.
\end{proof}

For centers $i<t$ the outermost fields are the earlier negative and the
later positive, up to commuting operators on a slice. If these choose
$\beta,\alpha$, sum their two slice differences first. Their factor is
\[
   \gamma^2\nabla B_i^\beta\,\Delta A_t^\alpha\,
                      Q_{-1,1}(Z_\alpha/Z_\beta).
\]
For the two middle fields replace $Z_\beta,Z_\alpha$ by
$\widetilde Z_\beta=pZ_\beta,\widetilde Z_\alpha=pZ_\alpha$.
Indeed the outer negative sends $1/Z_\beta$ to $p^{-1}/Z_\beta$;
the outer positive sends $Z_\alpha$ to $pZ_\alpha$.
Both selected roots are present at every middle cut on the appropriate
side. Their joint vacuum factor is the ratio $Q_{-1,1}$ above.
Successive normalized Cauchy factors multiply according to
$M(u,n)M(p^nu,m)=M(u,n+m)$. Applying this identity to each
horizontal--downward pair puts all remaining branch factors at the
two shifted roots, for either assignment of the middle fields.
In particular it also holds when another field selects the same root.
When the middle slices have their negative first, the two evaluations
agree. For the opposite order use Lemma~\ref{lem:quadraticinterchange}
with these shifted roots. Equal centers admit an order with both negative
fields first.

\begin{lemma}\label{lem:quadraticdomino}
Write $h=\beta+1$, $a=\beta+2$, and use the outer shifts at $\beta,a$
in $\Phi_{h,a}$. Then
\begin{multline}\label{eq:qdomino}
 l_1\langle C_2\rangle+
 \frac{l_1^2}{2}\big(\langle C_1^2\rangle-G_{11}\big)\\
 =l_1\gamma\sum_{\beta\,H,\,a\,D}
 B_{\beta+1}^\beta A_a^a\,Q_{-1,1}(Z_a/Z_\beta)\,\Phi_{h,a}.
\end{multline}
\end{lemma}
\begin{proof}
For each middle negative slice $j=t,t-1$ and positive slice $k=i,i+1$
the correction is supported at $k\le h<j$. In particular $\beta<i\le
h<t\le\alpha$. Write
\[
 U_k=\prod_{k\le r<h}R(Z_h/Z_r),\qquad
 V_j=\prod_{h<r<j}R(Z_h/Z_r),\qquad U_{h+1}=V_h=0 .
\]
At fixed $h,\beta,\alpha$ the left position sum is
$\sum_{i=\beta+1}^{h}(U_i-U_{i+1})\nabla B_i^\beta$.
By \eqref{eq:qtrans} the two transport parameters are $a'=Z_h$,
$b'=xZ_\beta$. If $a'\ne b'$ the identity
\begin{equation}\label{eq:qprimitive}
 (1-r_{a'})(1-sr_{b'})
       =\frac{(s-1)a'}{a'-b'}(r_{a'}-r_{b'})
\end{equation}
writes the summand as a forward difference of
$(s-1)a'U_i B_{i-1}^\beta/(a'-b')$, with both boundary values zero.
In the remaining case comparison of the $x$-degrees forces $h=\beta+1$,
and the sum is $B_h^\beta$.

For the right position sum separate $V_j\Phi_{h,j}-V_j$.
The subtracted part uses $B_h^\beta V_t=B_{t-1}^\beta$,
by \eqref{eq:qtrans}, including the extended zero at $t=h$.
Hence after taking the $t,t-1$ difference, multiplying by
$\Delta A_t^\alpha Q_{-1,1}$ and summing, this bare part is
$\langle C_2\rangle/\gamma$: the negative difference for $C_2$
occurs one slice earlier than for $C_1$.
For $W_t=V_t\Phi_{h,t}$ put $W_h=0$. Then
$W_{t-1}=r_{Z_{h+1}}(Z_t)W_t$ up to $t=\alpha$.
In fact, writing $z=Z_h/Z_{t-1}$, the forward quotient $W_t/W_{t-1}$
away from birth is $R(z)$ or $R(cz)$ for $h$ horizontal,
according as $t-1$ is horizontal or downward. For $h$ downward it is
$R(z/c)$ or $R(z)$, respectively. These are the asserted relations
since $Z_t/Z_{t-1}$ is $x$ or $x/c$; the cleared relation at birth has
zero numerator. All crossed steps lie strictly between the shifted
exterior roots. Applying \eqref{eq:qprimitive} a second time leaves only
$\alpha=h+1$ with value $A_\alpha^\alpha\Phi_{h,\alpha}$.
The chronological evaluation uses the negative of \eqref{eq:qmiddle}.
The two orders of distinct centers cancel the factorial $2$, and
$l_1\gamma\omega(1-s)=-1$. This gives \eqref{eq:qdomino}.
\end{proof}

\subsection{The two currents of each sign}
It remains to evaluate $(l_1^2G_{11}+l_2\langle\psi^2C_1\rangle)/2$.
Temporarily give the two selected downward steps distinct labels
$A=\{a,d\}$, their values $Z_\alpha$, and the two horizontal steps
distinct labels $B=\{b,e\}$, their values $xZ_\beta$. Remove the
opposite-sign $Q$ factors, which depend only on these four selections.
Remove also $(l_1\gamma)^2$ and the one-body weights taken at the left
cut for positives and the right cut for negatives. By
\eqref{eq:qtrans} the remaining one-body factors are products of the
$r_z(X)$ over the intervening slice weights $X$. For a set $J$ write
$r_J(X)=\prod_{z\in J}r_z(X)$.

Use $(q_+,V_+,m_+)=(p,s,1)$ and $(q_-,V_-,m_-)=(p^{-1},s^{-1},s)$.
For either sign define
\[
 K_\pm(t)=\frac{(1-t)(1-q_\pm t/V_\pm)}
                    {(1-q_\pm t)(1-t/V_\pm)},\qquad
 E_\pm(u,v)=\frac{V_\pm(u-v)(v-u)}
                  {(V_\pm u-v)(V_\pm v-u)} .
\]
Two distinct selections $u,v$ in chronological order supply
$K_\pm(v/u)$. Indeed dividing the second shifted singleton coefficient by
its unshifted value cancels all factors except that of the first
selected root; reverse the action for a suffix. The coefficient of the
elementary operator $e_2$ divided by the two singletons is $E_\pm$,
by multiplying their two mutual spectator factors and the elementary
prefactor $V_\pm$. Thus the coefficient for $p_2=e_1^2-2e_2$ is
$K_\pm(v/u)+K_\pm(u/v)-2E_\pm(u,v)$. The conversion of a second
power sum to two signed singleton factors contributes
$(V_\pm+1)/(V_\pm-1)$. The product of the two conversion factors times
$-s^2/(s^2-1)^2$ is precisely $\gamma^2$.

Here are the complete densities for one center of slice weight $X$.
For one pair put $d_X(a,b)=(1-r_a(X))(1-sr_b(X))$.
For $u,v$ of one sign write $k=K_\pm(v/u)$, $h=K_\pm(u/v)$,
$r_u=r_u(X)$, $r_v=r_v(X)$, $m=m_\pm$, and set
\begin{align*}
 T_\pm(u,v;X)
 &= (k+h)(1+m^2r_ur_v)-2m
 \begin{cases}r_u h+r_v k&(+),\\r_u k+r_v h&(-),\end{cases}\\
 H_\pm(u,v;X)&=(k+h-2E_\pm(u,v))(1-m^2r_ur_v).
\end{align*}
Indeed both fields at the native slice or both at the displaced slice
retain label order. With just one displaced field, that field is later
for the positive sign and earlier for the negative sign; there are two
labelings. The Adams field makes its slice choice only once.
For two complete pairs at $X$ the density is therefore
\[
 D_X(A,B)=\tfrac12 T_+(A;X)T_-(B;X)
             +\frac{l_2}{2l_1^2}H_+(A;X)H_-(B;X).
\]

\begin{lemma}\label{lem:twoslicetwo}
Through degree two the sum over assignments to centers is invariant
under permutations of the slice weights, after removing the exterior
one-body and opposite-sign factors.
\end{lemma}
\begin{proof}
Consider two consecutive weights $X,Y$. When both pairs are assigned
to these centers, their sum is
\begin{align}
 N(X,Y)={}&r_A(X)D_Y(A,B)+r_B(Y)D_X(A,B)\notag\\
 &+\sum_{\substack{u\in A\\v\in B}}
 K_+(u'/u)K_-(v'/v) r_v(Y)r_{u'}(X)
                   d_X(u,v)d_Y(u',v'),                      \label{eq:qtwoslice}
\end{align}
where primes denote complementary labels. The first term has both pairs
at $Y$, the second both at $X$, and the sum one at each. The factors
joining assignments before or after these two centers depend only on
their union. When one pair occupies the two-center window, its entry is
$r_b(Y)d_X(a,b)+r_a(X)d_Y(a,b)$, symmetric by \eqref{eq:qprimitive}
and continuation.

For two pairs we give the finite calculation of
$N(X,Y)-N(Y,X)$. Clear $\prod_{z\in A\cup B}(X-sz)(Y-sz)$.
The resulting polynomial has degree at most four in each slice
variable: each root occurs at most once per variable in any summand
of \eqref{eq:qtwoslice}. By simultaneous scaling take $A=(1,d)$,
$B=(b,e)$. In the field $\QQ(p,s,d,b,e)$ put
\[
       \frac{l_2}{l_1^2}=\frac{(s+p)(1-p)}{(s-p)(1+p)} .
\]
For $X=-i$, $Y=-j$, the numerators of the difference reduce to the
following entries, in increasing order of $j$,
\[
\begin{array}{c|cccc}
 i\ \backslash\ j&2&3&4&5\\ \hline
 1&0&0&0&0\\
 2&&0&0&0\\
 3&&&0&0\\
 4&&&&0
\end{array}
\]
One can reproduce the table by initializing the two terms on the first
line of \eqref{eq:qtwoslice}, adding its four terms for
$(u,v)=(1,b),(1,e),(d,b),(d,e)$, and subtracting the evaluation
with $i,j$ interchanged. All the denominators at these nodes are
nonzero rational functions. Antisymmetry gives the lower triangle
and the diagonal. Interpolation on five nodes in each variable proves
the identity; thus transposing any two adjacent slice weights preserves
the sum.
\end{proof}

We explain the specialization to repeated selections before permuting
the physical slice weights. If a downward step occurs twice, use the
two positive labels $pZ_\alpha,Z_\alpha$ in that chronological order.
For a repeated horizontal step use $xZ_\beta,pxZ_\beta$, in that
order. In the distinct-label identity divide by the ordered kernel for
the pair and specialize to this string. The kernel has a simple pole,
whereas its reverse and the elementary coefficient $E_\pm$ are
regular. A surviving pair must thus have the indicated order, with
weight one in place of its internal kernel. The divided coefficient
for $p_2$ is also one: a repeated shift occurs once in $e_1^2$ and
never in $e_2$. If both signs repeat, the two normalizing factors
involve independent variables. Consequently the specialization
reproduces the actual Macdonald action. At the second use of one
variable its one-body coefficient is taken at the shifted root,
omitting its own spectator in both coefficients. These weights propagate
as in \eqref{eq:qtrans} at the respective labels on the active cuts.
The common product of opposite-sign factors is legitimate also for
strings, since
\[
     Q_{-m,n}(w)=\prod_{i=0}^{m-1}\prod_{j=0}^{n-1}
                       Q_{-1,1}(p^{j-i}w).
\]

Take a window containing all the physical births $X=xZ_\beta=Z_{\beta+1}$
and deaths $X=Z_\alpha$. An assignment to center $i$ has exterior
transport $\prod_{k<i}r_a(X_k)$ for each positive and
$\prod_{k>i}r_b(X_k)$ for each negative, restricted to this window;
$D_{X_i}$ and $d_{X_i}$ include transport at the center itself.
For a repeated step the shifted copy may be continued beyond the
physical endpoint by its transport law. Every positive string uses
its unshifted copy last and every negative string uses that copy first
in actual slice order. An unshifted copy beyond its allowed range
therefore forces the zero $r_z(z)$, and the continuation introduces
no extra assignment. No physical transport denominator vanishes,
since $Z_i$ has $p$-degree zero and $sz$ has positive $p$-degree.
The specialized symmetry in Lemma~\ref{lem:twoslicetwo} also permits
coincident labels of opposite signs: its only label denominators
apart from transports involve labels of the same sign.

Suppose a selected downward step at $\alpha$ occurs more often than
the horizontal step at $\alpha-1$, counted with multiplicity zero if
absent. Move $X=Z_\alpha$ to the first center. Its unshifted positive
cannot occur later and thus all its copies must be assigned here.
A negative from any other step cannot be used here, since its
unshifted birth now occurs later. This center has insufficient negative
multiplicity. It follows that all nonzero assignments have matching
multiplicities at adjacent horizontal--downward pairs. Indeed an
excess at some downward step exists whenever the multiplicities fail
to match, since both totals are two. In the original slice order,
the first such turn must use all its positive copies there and receives
only its own negatives. The same argument at the next turn pins every
selected field to its corner.

\subsection{Identification of the height weights}
For a consecutive portion of the boundary starting with a horizontal
step $\beta$ and ending with a downward step $a$, let $w_i$ be the
box weights of the marked layers (including $p^{-1}$ for a second
layer) and set $z_i=c/w_i$. Relative to the weight on this portion
alone, the height polarization supplies the factor
\begin{equation}\label{eq:qexternal}
 \prod_i\left(
  \prod_{t<\beta,\,H}\frac{R(z_i/(xZ_t))}{R(z_i/Z_t)}
  \prod_{t>a,\,D}\frac{R(z_i/Z_t)}{R(cz_i/(xZ_t))}
 \right).
\end{equation}
To see this, the boundary term in the height polarization gives
$R(z_i/Z_t)$ at each horizontal--downward turn and its inverse at
each downward--horizontal turn, including the two terminal valleys.
Telescope along the horizontal runs to the left of the portion, and
along the downward runs to its right. The two-layer interaction
depends only on the ratio of the marked weights, so cancels in taking
this relative factor.

At the turn $\beta,a=\beta+1$ selected twice, let $A,B$ be the two
one-body residues there, and $\widetilde A,\widetilde B$ the values for
the second use of the same variable. Both the ordinary square and
the second power sum contribute
$\gamma^2 AB\widetilde A\widetilde B Q_{-2,2}(x)$ before their
respective symmetric-algebra scalars. Their exterior factors are
\eqref{eq:qexternal} for $z_1=Z_a$, $z_2=pZ_a$. On the one-box
diagram their common product reduces to
\[
 \frac{s^2Q_{-2,2}(x)}{R(x)R(px)R(c/x)R(pc/x)}=1,
\]
by $R(z)R(1/(sz))=s$ and the definition of $M$. The coefficient on
this portion is $(l_1^2+l_2)/2=(1-c)(1-pc)/((1-p)(1-p^2))$,
the weight of a height-two box.

For distinct turns write $t=Z_{\alpha_2}/Z_{\alpha_1}$ with the
earlier turn first. The product of their separate one-mark weights
must be multiplied by
\[
 K_+(t)K_-(t)Q_{-1,1}(xt)Q_{-1,1}(x/t)
     =R\big((p^{-1}-1)(1-x)(t+t^{-1})\big).
\]
Here $R$ is extended multiplicatively over a virtual character.
The identity follows by inserting the four fractions and reflecting
the factors in $t^{-1}$ with $R(z)R(1/(sz))=s$. The right side is
precisely the cross term of the height polarization. The two orders of
the marked pairs cancel the factorial in the ordinary square.
The second Adams operation cannot put a single insertion at two
different turns.

Finally the right side of \eqref{eq:qdomino} is supported on the words
$HHD$ and $HDD$. In both cases the two layer parameters in
\eqref{eq:qexternal} are $z_1=Z_{\beta+1}$, $z_2=Z_a$.
For every exterior horizontal $t$ the two residues
$B_{\beta+1}^\beta,A_a^a$ supply $R(Z_\beta/Z_t)/R(Z_a/Z_t)$.
For every exterior downward $t$ they supply $R(Z_a/Z_t)/R(cZ_\beta/Z_t)$.
The additional factor $\Phi_{\beta+1,a}$ multiplies the latter by
$S(Z_{\beta+1}/Z_t)$ on $HHD$, and the former by
$S(Z_{\beta+1}/(cZ_t))$ on $HDD$. Inserting the step increments gives
exactly \eqref{eq:qexternal}. The scalar on each isolated portion can
be seen directly in the following table. In the middle column $z=x$
for the horizontal domino and $z=c/x$ for the vertical one.
\[
\begin{array}{c|c|c}
 \text{portion}&l_1\gamma B_{\beta+1}^\beta A_a^a
       Q_{-1,1}(Z_a/Z_\beta)\Phi_{\beta+1,a}
       &\text{height weight}\\ \hline
 HHD&l_1(1-cz)/(1-pz)&l_1(1-cx)/(1-px)\\
 HDD&l_1(1-cz)/(1-pz)&l_1(1-c^2/x)/(1-pc/x)
\end{array}
\]
For example, set $Z_\beta=1$, $Z_{\beta+1}=x$ and
$Z_a=x^2$ or $x^2/c$, respectively. The horizontal residue at any cut
has own factor $1-s$, factors $R(Z_\beta/Z_t)$ for the other
horizontal steps before that cut and $1/R(cZ_\beta/Z_t)$ for its
downward steps. The downward residue has own factor $s-1$, factors
$1/R(Z_a/Z_t)$ on the horizontal prefix and $R(Z_a/Z_t)$ on the other
downward steps of the suffix. These products give the middle column.
For the last column the two-box row has removable box weight $x^{-1}$,
addable weights $x^{-2},x/c$, and layer sum $1+x^{-1}$; the column
has the same products with $x$ and $c/x$ interchanged, evaluated with
the corresponding two box weights in the given polarization.
Together these calculations account for every size-two reverse plane
partition and prove Theorem~\ref{thm:quadraticdisk}.

\section{The cubic coefficient}\label{sec:cubic}

\begin{theorem}\label{thm:cubicdisk}
For every Young diagram the two vacuum-normalized disk coefficients of
exchanged degree three satisfy $T_3=U_3$ in $\mathbb Q(p,c,x)$.
\end{theorem}

Write $M$ for the chronological moment, and $G$ for the moment obtained by
putting the negatives before the positives, preserving the order within
each sign. The argument begins by removing the scalar part of a change of order. We then prove the adjacent-slice permutation identity and compute the residues produced on restoring the original order.

\subsection{Removing the longer gaps}
Put $\mathcal C_h=l_1 C_h$ and $a_0=(1-s)^2(1-1/p)/(1-s/p)$, so
$l_1\gamma a_0=-1$. The scalar term in the difference of orders of a
positive at $k$ and a later negative at $j$ is
\[
 b(k,j)=a_0(1-s^{j-k})/(1-s)\quad (k<j),\qquad b(k,j)=0\quad(k\ge j).
\]
In a word of ordinary fields subtract this scalar at every inverted
pair by inclusion and exclusion. More precisely, set
\[
 M^\circ(w)=\sum_D(-1)^{|D|}\prod_{(k,j)\in D}b(k,j)\,M(w\setminus D),
\]
where $D$ ranges over matchings and $w\setminus D$ deletes both ends
of every matched pair. A single deletion in a signed sum joins
two intervals, since
\[
 b(i,j)-s b(i,j-1)-b(i+1,j)+s b(i+1,j-1)=a_0{\bf1}_{j=i+1}.
\]
The two untouched ends give the current of the joined interval. The
relative scalar is $l_1\gamma a_0=-1$. Consequently the six directed
single edges on three labeled insertions contribute $M(\mathcal C_1\mathcal C_2)$ to
$M^\circ(\mathcal C_1^3)/6$. The six directed paths of length two contribute
$M(\mathcal C_3)$. This proves that the whole third coefficient is
\begin{equation}\label{eq:cubic-moment}
 \tfrac16 M^\circ(\mathcal C_1^3)+\tfrac12 l_2 M(\mathcal C_1\psi^2 C_1)
                         +\tfrac13 l_3 M(\psi^3 C_1).
\end{equation}
Here $\psi^j$ is the $j$th Adams operation on the diagonal eigenvalues.
No subtraction is made in the second summand. The two fields of the
last summand are already in grouped order.

\subsection{Permuting the slice parameters}
We give explicitly the identity used to sum the slices. The two
alphabets $A,B$ below label positive and negative residues respectively.
The label of a positive variable of value $v$ is $v/s$, and that of
a negative variable of value $w$ is $xw$. Passing a slice parameter $X$
in the forward positive or backward negative direction multiplies
the corresponding one-body weight by
$r_z(X)=(X-z)/(X-sz)$. Initially take distinct indeterminates in each
alphabet. Put $r_J(X)=\prod_{z\in J}r_z(X)$.

For one sign let $(q,V)=(p,s)$ or $(1/p,1/s)$ as appropriate and write
\[
 K(t)=\frac{(1-t)(1-qt/V)}{(1-qt)(1-t/V)},\qquad
 E(J)=\prod_{\{z,w\}\subset J}
       \frac{-V(z-w)^2}{(Vz-w)(Vw-z)} .
\]
We also write $K_{q,V}$ to display the parameters, and
$K(I,J)=\prod_{z\in I,w\in J}K(w/z)$. Let $S(J)$ be the sum, over
the orders on $J$, of the products of the corresponding ordered pair
factors. For a power sum put
\[
 P(J)=
 \begin{cases}
 1,&|J|=1,\\
 S(J)-2E(J),&|J|=2,\\
 S(J)-3\sum_{u\in J}K(\{u\},J\setminus\{u\})E(J\setminus\{u\})+3E(J),
   &|J|=3 .
 \end{cases}
\]
These are Newton's formulas. The Cauchy determinant of an elementary
operator, divided by the singleton factors, is precisely $E(J)$
\cite[Props.~2.6--2.7]{Ko}.

For a partition $\nu=(j_1,\ldots,j_m)$ of $r$ in a fixed order, define
$T^\pm_\nu(A;X)$ as follows, using the appropriate sign throughout.
Sum over ordered decompositions $A=J_1\sqcup\cdots\sqcup J_m$ with
$|J_i|=j_i$, and over the two cuts for each group. Start with
$\prod_iP(J_i)$. The displaced cut of a positive group contributes
$-r_{J_i}(X)$, that of a negative group $-s^{j_i}r_{J_i}(X)$.
Multiply the pair factors $K(J_i,J_j)$ in cut order. Thus at one
center the positive at the center precedes a displaced positive, and
a displaced negative precedes the negative at the center. Use list
order at ties. For $|A|=|B|=r$ put
\begin{equation}\label{eq:cubic-density}
 D_X(A,B)=\sum_{\nu\vdash r} w_\nu T^+_\nu(A;X)T^-_\nu(B;X),
 \qquad w_\nu=\prod_{j\ge1}\frac{1}{m_j!}
              \left(\frac{l_j}{j l_1^j}\right)^{m_j},
\end{equation}
where $m_j$ is the multiplicity of $j$; the empty entry is one.
This is the contribution of a single slice after removal of its
one-body factors, the common opposite-sign contractions and
$(l_1\gamma)^r$. Indeed the scaled power sum has relative coefficient
$(V^j-1)/(V-1)^j$. The product for the two signs, times
$-s^j/(s^j-1)^2$, is $\gamma^j$.

For functions of paired, equally large alphabets let $*$ denote the
product obtained by summing over ordered decompositions
$(A,B)=(A_1,B_1)\sqcup(A_2,B_2)$ with $|A_i|=|B_i|$, inserting
$K_+(A_1,A_2)K_-(B_1,B_2)$ between the factors. The signs specify the
two choices of $(q,V)$. Write $R_X^+,R_X^-$ for multiplication by
$r_A(X),r_B(X)$, respectively. Pair factors make $*$ associative,
and $R_X^\pm$ preserve it.
\begin{lemma}\label{lem:cubic-slices}
In degrees at most three,
\begin{equation}\label{eq:cubic-commute}
 (R_Y^-D_X)*(R_X^+D_Y)=(R_X^-D_Y)*(R_Y^+D_X).
\end{equation}
Consequently the sum of the slice assignments is symmetric in their
parameters, after extracting the one-body weights at the two outer cuts.
\end{lemma}

\begin{proof}
The factors across the two slices are just the transports past $X$
for positives on the second slice and past $Y$ for negatives on the
first. Earlier centers precede later centers in the same-sign order,
including ties at the displaced cuts. The transports from outside the
pair depend only on its combined alphabets. Hence the second assertion
follows by adjacent transpositions.

Here are finite interpolation details for the new homogeneous part
of degree three. For each sign, index ten entries by pairs of partitions
$(\nu,\mu)$ with $|\nu|+|\mu|=3$. The positive entry $F^+_{XY}$ is
\[
 \sum_{I\sqcup J=A,\ |I|=|\nu|}
       r_J(X)K_+(I,J)T^+_\nu(I;X)T^+_\mu(J;Y).
\]
The negative entry $F^-_{XY}$ uses $B$, $K_-$, $T^-$ and $r_I(Y)$
in place of $r_J(X)$. Put $w_{\nu\mu}=w_\nu w_\mu$. The two sides of
\eqref{eq:cubic-commute} agree exactly when the rows
\[
 (F^+_{XY},-F^+_{YX}),\qquad
 (w_{\nu\mu}F^-_{XY},w_{\nu\mu}F^-_{YX})
\]
are orthogonal. The index $(\nu,\mu)$ runs through each of the two
ten-entry blocks. The following finite rational calculation suffices.
Take $Y=-1$, $X=-2,\ldots,-8$, $2\le p\le19$ and $21\le s\le50$,
with integral $X,p,s$. For each sign use every triple from the thirteen
primes from $11$ to $59$ as its labels. Form the two $286$ by $20$
rational matrices with the displayed rows. At each of the $3780$
choices of $(X,p,s)$ a row basis of the first matrix times the
transpose of the second is zero. The supplementary computation
implements these sums over subsets and uses rational row reduction.

For each unordered same-sign pair $z,w$ clear
$(z-pw)(w-pz)(z-sw)(w-sz)$. Each pair occurs in at most one ordered
or elementary factor of each summand. The resulting pair contribution
has degree at most two in each of $p,s$, four in each label. Clear also
$(X-sz)(Y-sz)$ for each label. The degrees are thus at most ten per
label and six per slice parameter. Negative displaced cuts supply at
most $s^3$. The only extra scalars are
\[
 \frac{l_2}{l_1^2}=\frac{(s+p)(1-p)}{(s-p)(1+p)},\qquad
 \frac{l_3}{l_1^3}=\frac{(s^2+sp+p^2)(1-p)^2}
                           {(s-p)^2(1+p+p^2)} .
\]
At most one occurs per summand. Clearing
$(s-p)^2(1+p)(1+p+p^2)$ gives final bounds $17$ in $p$ and $29$
in $s$. All grid denominators are nonzero. Symmetry in each alphabet
gives every ordering of its three distinct labels. With the other
two fixed, eleven nodes remain for each label; successively interpolating
therefore removes the distinct-node restriction and proves the identity
in the six labels. The seven slice values suffice for $X$.
Homogeneity under common rescaling restores $Y$, and the $18$ by
$30$ character grid finishes the proof. The lower degrees follow from Lemma~\ref{lem:twoslicetwo} and its degree-one calculation.
\end{proof}

\subsection{The change of order}
At a step $h$ write $\tau_h=Z_{h+1}/x$, which is $Z_h$ on an H step
and $Z_h/c$ on a D step. In the interchange of a positive at cut $k$
with a negative at $j>k$ a boundary residue is indexed by $k\le h<j$.
The two fields then have values $z_h=s Z_h$, $w_h=\tau_h$. Put
$S(v)=R(v)/R(cv)$ and
\[
 \begin{split}
 W_h(k,j)&=\prod_{k\le t<j,\ t\ne h} R(Z_h/Z_t),\\
 \Phi_{h,j}&=
 \begin{cases}
 \prod_{d\ge j,\ {\rm D}} S(Z_h/Z_d),&h\ {\rm H},\\
 \prod_{b<j,\ {\rm H}} S(Z_h/(cZ_b)),&h\ {\rm D}.
 \end{cases}
 \end{split}
\]
For chronological minus grouped order the residue has weight
$-a_0W_h(k,j)\Phi_{h,j}$. To check its action on the other fields,
expand the mixed kernel in the two orders. Its difference of expansions
is supported at $z=sw,pw$. At $sw$ the creation exponential of the
pair cancels and only contractions to its right remain; at $pw$ the
annihilation exponential cancels and only contractions to its left
remain. For a remaining field of sign $e=\pm1$ and value $u$, these
products are, respectively,
\[
 U_e(t)=K_{p,s^e}(t/s)K_{1/p,s^e}(t),\quad t=u/Z_h,\qquad
 V_e(t)=K_{p^e,s}(pt)K_{p^e,1/s}(t),\quad t=Z_h/(cu).
\]
The mixed kernel is unchanged as a rational function upon reversing
both the signs' order and the spectral ratio. Thus these same
products may be obtained by retaining all cross contractions in the
given same-sign order, even on the side where their product is one.
The two endpoints of the fused contour give $b(k,j)$ times the
moment with the pair deleted. The contour calculation may first be made
with independent boundary parameters, the horizontal ones sufficiently
large and the downward ones sufficiently small relative to all the
nested field circles. Its rational residues can then be specialized
to the $Z_h$ of the diagram.

\begin{lemma}\label{lem:cubic-poles}
After the scalar deletions in \eqref{eq:cubic-moment}, the strictly
inverted fields have either one or two boundary residues as above.
The remaining nonconstant choices are ordinary boundary residues and
their successive $p$-shifts, with one possible additional choice. At
an H residue it is a second negative with value $Z_h/c$; at a D
residue it is a second positive with value $sZ_h/c$. Either additional
choice occurs with just one inverted pair.
\end{lemma}

\begin{proof}
Consider first the ordinary cube. Number the fields of each sign in
their actual cut order, taking negatives first at ties. Every center
has its negative before its own positive, so each prefix contains at
least as many negatives as positives. The only candidate inversions are
$a=P_1N_2$, $d=P_2N_3$ and $b=P_1N_3$. Those among $a,d$ present
are adjacent disjoint pairs. Expand each as reversed order plus its
difference; after the two reversals only $b$ can require a swap.
There is one single difference for each inverted pair, and a double
difference precisely when both $a$ and $d$ occur. In the three possible
single differences the remaining words have orders
\[
 N_1[a]N_3P_2P_3,\qquad N_1N_2[b]P_2P_3,\qquad
 N_1N_2P_1[d]P_3 .
\]
In a double difference the intervals are disjoint since $N_2$
precedes $P_2$. The remaining fields are $N_1,P_3$.
In inclusion and exclusion, a matching using exactly one end of the
interchanged pair gives the same deleted word in both orders. The
matching of the pair itself subtracts its endpoint constant. Applying
this twice subtracts the two single constants and adds back their
product in the double difference. The resulting two boundary residues
retain their mutual contraction. For an earlier H and a later D it is
$U_+(pt)U_-(t)$, $t=Z_g/(cZ_h)$, where $g$ is the D step; otherwise
it is one. In the first case the outer H contour and inner D contour
move apart, so their mutual poles are never crossed.

In a single difference shrink the remaining positive contours from last
to first and enlarge the negatives from first to last. Their mutual
mixed poles stay on the other side. Repeating an ordinary boundary
selection gives its next $p$-shift by
\[
 R(t)K_{1/p,1/s}(t)=R(t/p)=R(t)K_{p,s}(1/t).
\]
At most one negative remains to the right of an H residue. At its
cut the one-body factor at $h$ and the two contractions with the pair
give, for $t=u/Z_h$,
\[
 R(t)U_-(t)=
 \frac{(1-t/s)(1-pt)(1-t/c)}
      {(1-ct)(1-t/p)(1-t)} .
\]
Thus the three choices at this scale have $t=1,p,1/c$. The other
negative in this word has cut at most $k$ and cannot also select $h$.
At a D residue at most one positive remains on the left. With
$t=u/(sZ_h)$ the product is
\[
 R(t)K_{p,s}(1/t)K_{p,1/s}(1/(pt)).
\]
It has the same three possible poles. The opposite orientations lie
outside the selected direction of contour motion. At distinct steps
the spectral monomials have distinct $x$-degrees.

The coincident choices $t=1$ cancel in pairs. On H the dressing of the
later negative at coincidence is $-1$, while the earlier negative is
undressed when the later one is fused. For the two eligible cuts $j_1,j_2$
the $j$-dependent part of the fused weight and the unfused H residue at
$Z_h$ have the same backward ratio $r_{x Z_h}(Z_j)$. Their products
agree upon exchanging $j_1,j_2$. All other factors have the same
evaluated variables and the same order for each sign. On D the
corresponding common forward ratio is $r_{Z_h}(Z_k)$ and the earlier
positive's dressing is again $-1$. This also gives cancellation at a
tie between two equally eligible fields. The endpoint weights of an
unfused ordinary positive and negative are $1$ and $s^j$ respectively.
Each contraction with the deleted endpoint tends to one. The primary
inequalities and the order of all the remaining fields are unchanged
when this field is deleted. In the sum over the candidate inversions
each such labeled inversion occurs once, so the two cuts of the
endpoint contribute $1-1$ or $s^i-s\,s^{i-1}$. For two inversions
the unfused fields are exterior and the same cancellation applies.

For the mixed summand apply $p_2=e_1^2-2e_2$ to the member of
the inverted pair coming from the second Adams operation.
Write $E_V(t)=-V(1-t)^2/((V-t)(Vt-1))$ for the two-variable
determinant factor, and put
\[
 S_{q,V}(t)=K_{q,V}(t)+K_{q,V}(1/t)-2E_V(t),\qquad
 \rho_V=(V^2-1)/(V-1)^2 .
\]
When the higher field is negative, each primary boundary residue is
multiplied by
\begin{equation}\label{eq:cubic-mixed}
 \rho_{1/s}\,\CT_u N_j(u)
 K_{p,1/s}(u/z_h)S_{1/p,1/s}(u/w_h).
\end{equation}
For a higher positive the multiplier is
$\rho_s\CT_u P_k(u)K_{1/p,s}(u/w_h)S_{p,s}(u/z_h)$.
The exterior negative and positive at the two centers may first be
evaluated on their prefix and suffix. Each chosen exterior root is
outside $[k,j)$, and its finite shift replaces that root by its
shifted value in the one-body factors. Equivalently, one retains
its regular contractions with the middle variables. The term with
no shift cancels in the corresponding signed difference.

Here is the contour order in \eqref{eq:cubic-mixed}. For two elementary
negatives take $0<|p|<|s|<1$ and equal radii for the two negative
variables. The determinant denominators $w_i-sw_j$ have their
prescribed expansions on these circles. Move the ordinary probe past
both. The two choices of the variable it meets cancel $1/2!$. At
$z=sw,pw$, with $t=u/w$ for the remaining variable, the respective
products $E_{1/s}(t)K_{p,1/s}(t/s)$ and $E_{1/s}(t)K_{p,1/s}(t/p)$
are
\[
 \frac{s(t-1)(pt-1)}{(st-1)(pt-s)},\qquad
 \frac{s(t-p)(t-1)}{(t-s)(st-p)} .
\]
The first has all its poles outside the unit circle, so $w$ can move
outwards, and the second has all its poles inside, so $w$ can move
inwards. Each gives one at the corresponding endpoint. For elementary
positives take $1<|s|<|p|$ and replace $t$ by $1/t$, moving the negative
probe. Both calculations continue rationally. For $e_1^2$ cross the
two ordinary fields successively. The two possibilities give
$K(t)+K(1/t)$ with the same mixed factor. Their two endpoint constants
cancel those of $2e_2$. This proves the stated contour formulas.

The function $S_{q,V}$ has a double zero at $1$ and is zero at $0$ and
$\infty$. At H set $t=u/Z_h$ in \eqref{eq:cubic-mixed}. The same-sign
denominators allow $p,1/p,s,1/s$. The mixed numerator cancels $1/p,s$,
and the one-body zero cancels $1/s$. The double zero at $1$ cancels
both the coincident poles. Besides ordinary H steps, the outward
contour thus has just $u=pZ_h,Z_h/c$. At D the extra scales stay on
the inside of this outward contour. The inward computation for a
higher positive interchanges H and D, giving $u=psZ_h,sZ_h/c$ at D;
its one-body zero is at $u=Z_h$. This exhausts the mixed term and
proves the lemma.
\end{proof}

\subsection{Summing the residues}
Record the step of each positive and each negative. An inverted pair
selects the same step for both signs. At an H step only this pair can
supply a positive, and at a D step only this pair can supply a negative.
For a choice with ordinary shifts, the assigned spectral values are
\[
  s p^t Z_h\quad\hbox{on the positive side},\qquad
  p^t \tau_h\quad\hbox{on the negative side},\qquad 0\le t<m
\]
at a step of multiplicity $m$ for the indicated sign. In the additional
choice of Lemma~\ref{lem:cubic-poles} replace $p$ by $1/c$ on the
single displaced value at its primary step.

\begin{lemma}\label{lem:cubic-pin}
In the signed sum over all cuts, a nonzero assignment must have as
many positives at $i$ as negatives at $i-1$ for every $i$. In this
case all the fields with these indices lie at their common center $i$.
\end{lemma}

\begin{proof}
We first identify the dependence on the cuts. At a primary step $h$
split $W_h(k,j)=U_k^h V_j^h$ into the products over $k\le t<h$ and
$h<t<j$, extending $U_{h+1}^h=V_h^h=0$. Give the positive the factor
$U_k^h$ and the negative $V_j^h\Phi_{h,j}$. The two transport ratios are
\[
 U_{i+1}^h/U_i^h=r_{Z_h}(Z_i),\qquad
 (V_{i-1}^h\Phi_{h,i-1})/(V_i^h\Phi_{h,i})=r_{Z_{h+1}}(Z_i).
\]
For the second ratio a horizontal primary gives forward factors
$R(v)$ or $R(v)/S(v)$ through H or D, with $v=Z_h/Z_{i-1}$.
A downward primary gives $R(v/c)$ or $R(v)$. Substitution of the
two possible steps proves the identity, including its cleared endpoint.
For all other one-body factors the same relations follow by dividing
$P_{i+1}/P_i$ or $N_{i-1}/N_i$ at the assigned spectral value.

The cut-independent factors are the one-body weights at the two outer
cuts, the primary residue constants and the mixed kernels at their
assigned values (with the primary kernel omitted). What remains is
\eqref{eq:cubic-density} with its transports, applied to the labels
$p^t Z_h$ and $p^t Z_{h+1}$ of either sign. To see that all allocations
occur with their stated coefficients, temporarily distinguish all
slots of each sign. A fixed primary at $h$ requires exactly the strict
inversion $k\le h<j$. Once its two slots are assigned, these endpoint
conditions are independent. For ordinary fields each allowed single
pair occurs once in the three inversions of the preceding proof. Two
disjoint inversions necessarily use $a,d$, giving the single double
difference. Their four mutual contractions have just the product of
the two same-sign and two mixed kernels, as in that proof. For an
inverted second power sum the two possible singleton choices and the
simultaneous elementary choice have precisely the coefficients of
$P(J)$ by \eqref{eq:cubic-mixed}. An allocation in a single bilinear
insertion cannot straddle a primary step, since its negative cut is
at most its positive cut. Thus including all three partitions of
three in the density adds no spurious primary terms. The factorials
of indistinguishable insertions are the $m_j!$ in
\eqref{eq:cubic-density} after summing the labeled allocations.
These arguments also cover the grouped term, which has no primary.
An ordinary unshifted cutoff is enforced by its zero transport at the
physical end of its interval.

For an ordinary repeated selection, specialize the labels of a length
$m$ string to $z,pz,\ldots,p^{m-1}z$. First divide the distinct-label
density by the ordered kernels of the consecutive copies, and take
the limit with independent consecutive ratios. There is one simple
pole on each such divisor for decreasing levels in positive cut order
and increasing levels in negative cut order. Reverse order and
membership in a single elementary subset are regular. Nonconsecutive
copies in a common subset cannot survive either, because the missing
intermediate copy must be strictly between them. All other same-sign
factors are regular, including the pair at distance two. Hence the
limit keeps exactly the ordered $p$-shifts; the difference between the
divided kernel and its logarithmic residue is a cut-independent
scalar. For a step carrying a primary, at most one sign is repeated.
Its upper mixed kernel is regular for an ordinary $p$-shift.
Equivalently this procedure reproduces successive singleton shifts of
one Macdonald variable, while an elementary subset shifts distinct
variables only. The same-sign orders are unaffected by the inserted
mixed kernels. The physical support of each upper slot is automatic,
since its unshifted positive is used last, or its unshifted negative first.

For the additional choice at H or D no same-sign limiting operation is
needed. The mixed kernel at that pole is replaced by its logarithmic
residue, $(1-1/p)(1-s)/(1-s/p)$ at H and its negative at D. The
opposite same-sign order has zero factor
$K_{1/p,1/s}(c)=0$ at H or $K_{p,s}(1/c)=0$ at D. Consequently
the base negative precedes the additional negative at H, and the
additional positive precedes the base positive at D, unless both
are in one simultaneous elementary subset. In the latter case the
two share their cut and retain $E(J)$. The label of the additional
slot is the label of its base divided by $c$. Transporting it throughout
the interval adds no choices: its base has the genuine endpoint zero
and the opposite strict order vanishes. The formula is therefore the
same distinct-label density, specialized at regular same-sign
arguments, with the same endpoint constraints.

Now extract the initial positive and final negative one-body weights
in a finite window containing all selected steps and their successor
vertices. In this window the cut sum is invariant under reordering
the parameters $Z_i$ by Lemma~\ref{lem:cubic-slices}. This remains
true after the repeated-label limits. Transport denominators cannot
vanish at a physical slice: if $Z_i$ and the base label have the
same $x$-degree, their ratio is one, before the specified powers of
$p$ or $1/c$ are inserted. Suppose more positives occur at $i$ than
negatives at $i-1$, and bring $Z_i$ to the first slice. Its unshifted
positive cannot pass this slice because $r_{Z_i}(Z_i)=0$. All its
upper positives must occur there as well. A negative based at another
vertex cannot be used here, since the backward transport of its base
would cross its own zero and the base must occur first. There are too
few available negatives for equal degrees of both signs on this slice.
If any multiplicity disagrees there is such an excess somewhere.
When all multiplicities agree, proceed in physical order. At the first
occupied vertex all its positives must have been used by this cut
and only its own negatives can yet be used. Their common center and
native cut are forced. Removing these slots and proceeding to the
next vertex proves the assertion.
\end{proof}

The ordinary matched residues have a simple interpretation. Consecutive
occupied contents are joined across their shared primary step. An
occupied interval opens at H and closes at D and consists of the
southeast-most box on each intervening content. A multiplicity greater
than one requires an H birth and a D death, hence is possible only at
a removable corner. Conversely, at size three each layer contains
at most one box on each diagonal: two on the same diagonal force a
square of four. The connected components of the bottom layer are
therefore border strips with successive contents, and an upper box
at a nonremovable corner would force its neighbor to have height
at least two, requiring at least four boxes in total. This gives the bijection between
the ordinary matched patterns and the reverse plane partitions of
size three.

\subsection{Weights of the surviving heights}
We compare their full weights. For a layer at content $i$ and height
$t$ counted from zero put $z=p^tZ_i$. Its positive and negative
variables have values $sz,z/x$, even when one of them was fixed
by fusion with a neighboring layer. For two such layers in increasing
content order their four regular cross contractions have product
\begin{equation}\label{eq:cubic-cross}
 \begin{split}
 \mathcal K(y)={}&K_{1/p,1/s}(y)K_{p,s}(y)
          K_{1/p,s}(sxy)K_{p,1/s}(y/(sx))\\
       ={}&R\big((p^{-1}-1)(1-x)(y+y^{-1})\big),\qquad
             y=z_{\rm later}/z_{\rm earlier}.
 \end{split}
\end{equation}
In the last expression $R$ acts multiplicatively on a character. The
identity follows by expanding the four binomial quotients. It gives
exactly the cross term of the quadratic height polarization between
the two layers, whose characters are $c/z_{\rm earlier}$ and
$c/z_{\rm later}$.

At a single corner of content $i$ a stack of $k\le3$ copies can use
only $e_1^j$ in each Newton power, since an elementary subset cannot
shift the same root twice. Its scalar is
\[
 h_k[l]=[v^k]\exp\!\sum_{j\ge1} l_j v^j/j
       =\prod_{t=1}^k(1-cp^{t-1})/(1-p^t).
\]
The last equality follows by comparing the generating function at $v$
and $pv$. The selected H and D variables give opposite finite Cauchy
shifts, which cancel their mutual bridge. Each layer $z=p^tZ_i$ has
the exterior factor
\begin{equation}\label{eq:cubic-outside}
 \prod_{b<i-1,\ {\rm H}}\frac{R(z/(xZ_b))}{R(z/Z_b)}
 \prod_{d>i,\ {\rm D}}\frac{R(z/Z_d)}{R(cz/(xZ_d))}.
\end{equation}
Indeed the horizontal Macdonald and downward Cauchy singleton
factors at an earlier $b$ give $s^{-1} R(z/(xZ_b))/R(z/Z_b)$,
and the two factors at a later $d$ give
$sR(z/Z_d)/R(cz/(xZ_d))$. The counts of the two types cancel
the negative normalization $s^i$ for each layer. On the one-box
diagram the height weighs $s^{-k}\prod_{t=1}^k R(p^{-t})=h_k[l]$.
On a general diagram each exterior step contributes to the linear
part of the polarization separately for each layer, giving
\eqref{eq:cubic-outside}. This proves the equality for a single stack;
\eqref{eq:cubic-cross} now treats disjoint stacks and interactions
between separate border strips.

For completeness we compute the one-layer weight of a connected
three-box strip. Write $\beta,\alpha$ for its opening H and closing
D, $b=Z_\beta$, $d=Z_\alpha$. Its internal H and D step weights are
denoted by $u$ and $v$, respectively. Delete the strip, leaving
$\mu\subset\lambda$. If $Q$ is the sum of its box characters and
$E_\mu$ the sum of those of $\mu$, reflection
$R(y)R(1/(sy))=s$ in the polarization gives its weight
\[
 R\big(\big[(1-x)(1-c/x)E_\mu-c\big]Q^*\big).
\]
In fact the unreflected increment is
$Q/p+(1-x)(E_\mu Q^*-QE_\mu^*/p)$; the scalar from
reflecting the first term cancels $s^{-|\,\lambda/\mu\,|}$.
Each step outside $[\beta,\alpha]$ gives
\eqref{eq:cubic-outside}, with $e<\beta$ for the H steps and $f>\alpha$ for the D steps in its
corresponding product, for each of the three layer parameters $z$.
On the residue side the opening and closing singletons furnish
$R(Z_\beta/Z_e)/R(Z_\alpha/Z_e)$ at an earlier H step $e$ and
$R(Z_\alpha/Z_f)/R(cZ_\beta/Z_f)$ at a later D step $f$.
Multiplication by the intervening $\Phi$ factors gives the same
exterior products by the two step recursions for $Z_i$.

It remains to compare on the isolated strip. The two endpoints and
their mutual kernel reduce to the scalar $l_1$, since
$l_1\gamma a_0=-1$. Every internal H leaves $1/R(cd/u)$ and every
internal D leaves $1/R(c^2 b/v)$ after cancellation against its
endpoint contractions. For an H preceding a D their two factors
from $\Phi$ and the mutual contraction $U_+(pt)U_-(t)$,
$t=v/(cu)$, have product $S(v/u)S(cu/v)$. These identities follow
on substituting the displayed $U_e$ and $S$. Thus the character
whose image under $R$ gives the residue is
\begin{equation}\label{eq:cubic-ribbon-patch}
 -c-\sum_u cd/u-\sum_v c^2b/v+
       (1-c)\sum_{u\ {\rm before}\ v}(v/u+cu/v).
\end{equation}
One can check the height expression on the four possible patches
directly. Put $Y=d+\sum u+\sum v$. Removing the strip divides the
intervening vertex weights by $c$ and exchanges the initial H and
final D. Telescoping $(1-x)(1-c/x)E_\mu-c$ along the horizontal
steps of the remaining patch gives the height character
$c\big((x^{-1}-1)\sum u^{-1}-1/d\big)Y$.
Substitution of the two step ratios gives the following common
values for this character and \eqref{eq:cubic-ribbon-patch}:
\[
\begin{array}{c|l}
 {\rm HH}&-c(1+x+x^2)\\
 {\rm HD}&-c-x^2-c^2/x^2+(1-c)(x+c/x)\\
 {\rm DH}&-c-cx-c^2/x\\
 {\rm DD}&-c(1+c/x+c^2/x^2).
\end{array}
\]
This also exhibits the mutual factor on the strip with internal
word HD. The analogous one-internal-letter calculation is the
domino of the preceding section.

There remains the domino with a doubled corner. At this corner
the first negative and the last positive are the base fields. Keep
the other two variables independent while taking the residue of the
base domino. Taking their same-sign residues next shifts each
spectrum by $p$, as in the single stack. The mixed kernel with
the neighbor stays regular: its later-to-earlier ratio in
$K_{p,1/s}$ is $p/s$ on HHD and $1/p^2$ on HDD. The local
factor at the corner is thus the ratio $W_2/W_1$ of the stack
and singleton weights already calculated. The upper layer still
interacts with the neighbor, giving \eqref{eq:cubic-cross}. Hence
the result in the two cases is
\[
 W_{\rm domino}\frac{W_2}{W_1}\,\mathcal K(y),\qquad
 y=px\quad({\rm HHD}),\qquad y=x/(cp)\quad({\rm HDD}).
\]
The quadratic height character gives exactly this expression by
adding the upper layer. Together these calculations account for
every ordinary matched pattern.

\begin{proof}[Completion of the proof of Theorem~\ref{thm:cubicdisk}]
For an additional pole of Lemma~\ref{lem:cubic-poles},
Lemma~\ref{lem:cubic-pin} leaves only a doubled removable end of a
domino. At that center one sign has an ordinary $p$-string. On the
other sign take $(q,V,t)=(1/p,1/s,1/c)$ on H and $(p,s,c)$ on D,
where $t$ is the ratio of the later to the earlier value in the
permitted order. For two ordinary insertions the unique order on
the $p$-string and the zero of the reverse-order kernel on the other
sign leave just $l_1^2 K_{q,V}(t)/2$. For one second power sum the
ordinary string again allows just its singleton order, while the
other sign contributes its two-variable Newton factor. The sum,
after removal of their common factors, is
\[
 \tfrac12\big[(l_1^2+l_2)K_{q,V}(t)-2l_2 E_V(t)\big]=0 .
\]
Indeed in both cases $K_{q,V}(t)/E_V(t)=
2l_2/(l_1^2+l_2)$ by substitution. This removes all the additional
residues. The ordinary matched residues give the complete height sum,
which proves the theorem.
\end{proof}

For example on the two-box row the contribution of the additional
pole in the mixed summand, including its coefficient in
\eqref{eq:cubic-moment}, is
\[
 \frac{c(c+1)p(x+1)(x-1)^2(c-1)^2(x-c^2)}
 {2(1-p)(x-pc)(px-c)(pc-1)(px^2-c)} .
\]
At $c=3,x=5$ its residue and constant coefficient in $1-p$ are
$-576/11$ and $-16704/121$. Both are canceled by the ordinary cube.

\section{Descent to the Picard variety}\label{sec:family}

Let $C$ be a smooth projective complex curve of genus $g=b+1\geq1$,
$B=\operatorname{Pic}^{b}(C)$, and $L$ the Poincar\'e line on
$C\times B$, trivialized over a chosen point of $C$.
Use the ordered determinant lines of \cite[Lemma~2.1]{W1} for
the two stable-pair indices. Write $q$ for the original Euler
variable, $t$ for the exchanged grading, $c=\kappa^{-1}$ and
$\sigma^2=\kappa$. For a partition $\lambda$ of $m$, set
$H_\lambda=\sum_{v\in\lambda}h_v$, where $h_v$ and $\ell_v$
denote the hook length and leg length. The scalar removed from each
of its two contributions in passing to the disk calculation is
\begin{equation}\label{eq:family-scalar}
 (-1)^{mb}q^{-bH_\lambda}
 c^{\,b\sum_{v\in\lambda}(\ell_v+1/2)}.
\end{equation}
Here $c^{1/2}=\sigma^{-1}$. Thus disk degree $(n,k)$ corresponds
to $q^{-bH_\lambda+n}t^k$ on the curve.

There is a second normalization in the current calculation. With
$p$ the disk cotangent character and $s=p/c$, the common series of
exchanged degree zero is
\begin{equation}\label{eq:family-axis}
 A_{\lambda;c,s}(x,0)= A_{\lambda;s,c}(0,x)
 =\prod_{v\in\lambda}F_s(x^{h_v}c^{-\ell_v}),\qquad
 F_s(z)=\sum_{i\geq0}\frac{(s;p)_i}{(p;p)_i}z^i .
\end{equation}
Here $(a;p)_i=\prod_{j=0}^{i-1}(1-ap^j)$.
The hook product is \cite[Lemma~5.1]{W1}. Its constant term is one.
Consequently an identity of rational functions in $p,c,x$ for a
coefficient divided by \eqref{eq:family-axis} gives equality of the
full coefficient as a power series in $x$, on multiplying by
\eqref{eq:family-axis}.

\begin{proposition}\label{prop:family}
Fix $d\geq0$ and a partition $\lambda$ of $m$. Suppose the two disk
series $A_{\lambda;c,s}(x,y)$ and $A_{\lambda;s,c}(y,x)$ agree
through $y$-degree $d$ in $\mathbb Q(p,c)[[x,y]]$. Then the
contributions of $\lambda$ to the two family indices agree in
\[
 K^0(B)[\sigma^{\pm1},(1-\sigma^j)^{-1}\ (j\geq1)]
\]
at every $q^{-bH_\lambda+n}t^k$ with $\min(n,k)\leq d$ and
$n,k\geq0$. If the hypothesis holds for every $\lambda\vdash m$,
the two degree-$m$ indices, multiplied by $\sigma^{mb}$, agree
modulo $t^{d+1}$ in
\[
 q^{-bm(m+1)/2}\,K^0(B)[\kappa^{\pm1}][[q,t]].
\]
Here $K^0(B)$ denotes integral algebraic $K$-theory.
\end{proposition}

\begin{proof}
The involution $(x,y,c)\mapsto(y,x,p/c)$ gives the disk identity
also through $x$-degree $d$. For any $(n,k)$ in either strip it
therefore gives every coefficient of the rectangle
$0\leq a\leq n$, $0\leq e\leq k$.
We recall how the two-divisor descent of
\cite[Proposition~7.1]{W1} uses these coefficients.
After removal of \eqref{eq:family-scalar}, the nested divisors of the
fixed boxes and the insertions from the moving factors give classes
on $C^{(n)}\times C^{(k)}\times B$. On disjoint supports their
restrictions are exterior products, with the graded permutation
action on equal pairs of multiplicities.

In an affine coordinate, the derived restriction at a single point
of multiplicities $(a,e)$ is the disk coefficient times
\[
    z^e(p;p)_a(p;p)_e ,
\]
where $z$ is the character of the frame of $L$. The two finite
products are the cotangent Euler classes of the symmetric powers
of the affine line at their origins. Hence the comparison is made
while $p$ and $z$ are independent characters. Expanding the
zero-counting-degree Euler inverses in the volume character, every
coefficient of the restriction is an integral torus character.
It follows that the two restrictions agree integrally. On a cycle
of a permutation of equal support types the supertrace is the
punctual supercharacter evaluated on the product of the torus
entries around that cycle. This includes the cycles of lengths two
and three at the corresponding exchanged degrees.

The derived restrictions detect the two classes on the stack of
coordinate jets and frames used in \cite[\S7.2]{W1}. Indeed,
filter by number of distinct support points. The conormal Euler
of each stratum in its step of the filtration has constant term
one and all other terms of positive simultaneous dilation weight.
It acts injectively on the integral representation group of the
stabilizer, including permutations of equal types. Coherent
localization and the self-intersection formula then give detection
by induction. Pullback to the two symmetric powers uses the actual
cotangent and Poincar\'e lines of the curve. Proper pushforward to
$B$ and restoration of \eqref{eq:family-scalar} give the asserted
diagram equality.

Finally, \cite[Lemma~2.2]{W1} puts each normalized degree sum in
the indicated Laurent polynomial ring coefficientwise. At fixed
$q^j t^k$ with $k\leq d$ a diagram occurs at excess
$n=j+bH_\lambda$ if this is nonnegative. Each such diagram has
been compared. Localization at $1-\sigma^i$, $i>0$, is injective
on $K^0(B)[\sigma^{\pm1}]$, by its extremal unit coefficient,
and $\kappa\mapsto\sigma^2$ is injective there as well. The
summed equality thus holds with the stated integral coefficients.
\end{proof}

\begin{thebibliography}{Mon26}
\bibitem[BC14]{BC} A.~Borodin and I.~Corwin, \emph{Macdonald processes}, Probab. Theory Related Fields \textbf{158} (2014), 225--400. \href{https://doi.org/10.1007/s00440-013-0482-3}{doi:10.1007/s00440-013-0482-3}.
\bibitem[DS22]{DS} H.~Dinkins and A.~Smirnov, \emph{Quasimaps to zero-dimensional $A_\infty$-quiver varieties}, Int. Math. Res. Not. IMRN (2022), 1123--1153. \href{https://doi.org/10.1093/imrn/rnaa129}{doi:10.1093/imrn/rnaa129}.
\bibitem[Ko21]{Ko} S.~Koshida, \emph{Free field theory and observables of periodic Macdonald processes}, J.~Combin. Theory Ser.~A \textbf{182} (2021), 105473. \href{https://doi.org/10.1016/j.jcta.2021.105473}{doi:10.1016/j.jcta.2021.105473}.
\bibitem[Mon26]{Mon} S.~Monavari, \emph{The refined local Donaldson--Thomas theory of curves}, Geom. Topol. \textbf{30} (2026), 3423--3484. \href{https://doi.org/10.2140/gt.2026.30.3423}{doi:10.2140/gt.2026.30.3423}.
\bibitem[NO16]{NO} N.~Nekrasov and A.~Okounkov, \emph{Membranes and sheaves}, Algebr. Geom. \textbf{3} (2016), 320--369. \href{https://doi.org/10.14231/AG-2016-015}{doi:10.14231/AG-2016-015}.
\bibitem[W26a]{W0} B.~J.~Wuebben, \emph{The topological string on balanced local curves from two circle actions on a Calabi--Yau fivefold}, version of November 26, 2026. \href{https://doi.org/10.5281/zenodo.24549724}{doi:10.5281/zenodo.24549724}.
\bibitem[W26b]{W1} B.~J.~Wuebben, \emph{The first variation and the one-mark coefficient of circle exchange for balanced local curves}, November 26, 2026. \href{https://doi.org/10.5281/zenodo.24570158}{doi:10.5281/zenodo.24570158}.
\end{thebibliography}

\end{document}
